\section{Method}
\label{sec:method}

\subsection{Direction B: Representation-Space Reassignment (Intervention/Transport)}
\label{sec:dirB}

Direction B performs reassignment in a \emph{concept representation space} extracted from internal activations.
Let $\Phi_{\theta}(v)\in\mathbb{R}^{d}$ be a concept prototype embedding for node $v$ under model $\theta$.

\paragraph{Choice of concept representation $\Phi$.}
We consider two practical instantiations:
\begin{itemize}
\item \textbf{Activation prototype:} $\Phi_{\theta}(v)$ is the mean of a fixed internal activation vector (e.g., a cross-attention context/value vector) collected under prompts $p\sim\mathcal{P}_v$.
\item \textbf{SAE code prototype (optional):} train an SAE on the same activations, and let $\Phi_{\theta}(v)$ be the mean sparse code.
\end{itemize}
Both are compatible with Stable Diffusion; the SAE variant often yields a cleaner, more interpretable geometry.

\paragraph{Target prototype via probabilistic anchors.}
Using prototypes from the \emph{original} model,
\begin{equation}
r_v \triangleq \Phi_{\theta_0}(v),
\qquad
r_u^{\star} \triangleq \sum_{a\in\mathrm{A}(u)} \pi_{u\rightarrow a}\, r_a.
\label{eq:target-proto}
\end{equation}

\paragraph{Intervention operator.}
Let $y$ denote the internal activation vector used by $\Phi$ (or its SAE code).
We learn a \emph{transport} operator $T(\pi)$ that depends on the anchor distribution:
\begin{equation}
\tilde{y} = T(\pi_{u\rightarrow \cdot})\, y,
\qquad
T(\pi) \triangleq \sum_{a\in\mathrm{A}(u)} \pi_{u\rightarrow a}\, T_a,
\label{eq:transport}
\end{equation}
where each $T_a$ is a small trainable linear map (optionally low-rank) acting on the activation space.

\paragraph{Prototype matching loss.}
Define the post-intervention prototype for $u$ as
\begin{equation}
\Phi_{\theta_0,T}(u) \triangleq \mathbb{E}_{x_0,t,\epsilon}\ \mathbb{E}_{p\sim\mathcal{P}_u}\big[\tilde{y}\big].
\end{equation}
We train $T$ (and the GCN producing $\pi$) to match the target prototype:
\begin{equation}
\mathcal{L}_{\mathrm{proto}}(u) =
\big\| \Phi_{\theta_0,T}(u) - r_u^{\star} \big\|_2^2.
\label{eq:proto-loss}
\end{equation}

\paragraph{Preservation loss.}
For retain concepts (or retain prompts) we constrain interventions to be identity:
\begin{equation}
\mathcal{L}_{\mathrm{pres}} =
\mathbb{E}_{p\sim\mathcal{P}_{\mathrm{ret}}}\big[\|\tilde{y}-y\|_2^2\big].
\label{eq:pres-loss}
\end{equation}

\paragraph{Final objective (Direction B).}
\begin{equation}
\min_{T,\,\mathrm{GCN}}
\ \ \mathcal{L}_{\mathrm{proto}}(\cdot)
+ \lambda\,\mathcal{L}_{\mathrm{pres}}
- \beta\,\mathcal{R}_{\mathrm{ent}}(\cdot)
+ \eta\,\sum_{a}\|T_a\|_F^2.
\label{eq:objective-B}
\end{equation}

\subsection{Concept Space Analysis}
\label{sec:space}

To interpret unlearning as geometric deformation, we analyze prototypes $\{r_v\}_{v\in V}$ before and after unlearning/intervention.
We report (i) reassignment error $\|r_u' - r_u^{\star}\|$, (ii) neighborhood distortion for retain regions,
\begin{equation}
\Delta_{\mathrm{retain}} \triangleq
\mathbb{E}_{(v,w)\in \mathcal{E}_{\mathrm{retain}}}
\big|d(r_v',r_w') - d(r_v,r_w)\big|,
\end{equation}
and (iii) visualizations (PCA/UMAP) of trajectories $r_u \rightarrow r_u'$ relative to anchor barycenters.

\begin{algorithm}[t]
\caption{Probabilistic Concept Reassignment via Graph-Guided Anchors}
\label{alg:main}
\begin{algorithmic}[1]
\REQUIRE Graph $G=(V,E)$, target concept $u$, prompt sets $\{\mathcal{P}_v\}_{v\in V}$, retain prompts $\mathcal{P}_{\mathrm{ret}}$
\STATE Compute node embeddings $\{z_v\}$ using GCN (Sec.~\ref{sec:gcn})
\STATE Construct anchor candidates $\mathrm{A}(u)$ using Eq.~\eqref{eq:anchor-set-leaf} or Eq.~\eqref{eq:anchor-set-nonleaf}
\STATE Compute reassignment distribution $\pi_{u\rightarrow a}$ using Eq.~\eqref{eq:pi}
\IF{Direction A (LoRA/ESD)}
    \STATE Form mixture teacher $\epsilon_{\mathrm{mix}}$ by Eq.~\eqref{eq:eps-mix}
    \STATE Update $\theta'$ by minimizing Eq.~\eqref{eq:objective-A}
\ELSE
    \STATE Compute prototypes $\{r_v\}$ and target barycenter $r_u^\star$ by Eq.~\eqref{eq:target-proto}
    \STATE Learn transport $T(\pi)$ by minimizing Eq.~\eqref{eq:objective-B}
\ENDIF
\end{algorithmic}
\end{algorithm}

